\documentclass[12pt]{article}
\usepackage[utf8]{inputenc}
\usepackage[T1]{fontenc}
\usepackage{hyperref}
\title{Pazeh Language Resources for K-12 Curriculum Development}
\author{Compiled by Claude Sonnet 5.5 (claude-sonnet-5-5) for Andr\'as Kornai}
\date{}
\begin{document}
\maketitle

\section{Introduction}
Pazeh (ISO 639-3 pzh) is an Austronesian language of central Taiwan. Pazeh and Kaxabu are generally treated as two closely related varieties. A 2007 Paradisec blog post reported a single full speaker (aged 94) in the Pazeh community, where language lessons for children had begun, while Kaxabu elders in Shoucheng were doing revitalisation work. That figure is dated and current numbers may differ.

\section{Primary Educational Resources}
No printed primers or textbooks for Pazeh were found in the searches. The Paradisec post mentions children's lessons and Kaxabu stories and songs, but no published materials are cited there.

\section{Linguistic and Reference Materials}
\begin{itemize}
\item Lin Chihkai, corpus-based study of Pazeh and Kaxabu (NINJAL abstract, PDF), which draws on Li and Tsuchida (2001), \emph{Pazih Dictionary} (Academia Sinica), for Pazeh data: \url{https://www2.ninjal.ac.jp/past-events/Methodsxvi/abs/Lin_Chihkai.pdf}
\end{itemize}

\section{Religious and Community Resources}
No scripture translations or community organisations were identified in the search results.

\section{Digital and Audio Resources}
\begin{itemize}
\item Paradisec blog, ``Endangered languages and Taiwan'' (field visit to Pazeh and Kaxabu communities; Paradisec is an archive of endangered-language recordings): \url{https://www.paradisec.org.au/blog/2007/06/endangered-languages-and-taiwan/}
\end{itemize}

\section{Curriculum Development Context}
Pazeh has no known school teaching and almost no first-language speakers. A curriculum would be a revitalisation effort built on the 2001 dictionary and archival materials, with the Kaxabu revitalisation work as a possible partner.

\end{document}